\chapter{Activities of the Department}

During the internship, I was part of the OMS (Operations Management System) team at Hyperface and worked in the role of a Product Management Intern. The OMS team is responsible for managing and improving internal operational workflows, transaction monitoring systems, issue tracking processes, and platform coordination activities that support the company’s fintech products and services. The team works closely with engineering, quality assurance, operations, and business stakeholders to ensure smooth execution of product features and operational processes.

As a Product Management Intern, I assisted in understanding product requirements, analyzing workflows, documenting features, coordinating with development teams, and supporting product planning activities. The role involved collaborating with different teams to improve operational efficiency, monitor product functionality, and contribute to the development of scalable fintech solutions. The department follows a collaborative and fast-paced work environment focused on innovation, problem-solving, and continuous product improvement.

\section{Departmental Role and Objectives}

The OMS (Operations Management System) team at Hyperface is responsible for managing and improving operational workflows, transaction handling systems, issue tracking processes, and internal platform coordination activities. The department focuses on ensuring that fintech products and digital credit solutions operate efficiently, securely, and reliably. As part of the product management function, the team works closely with engineering, operations, quality assurance, and business teams to streamline processes and support smooth product execution.

The OMS team contributes significantly to the company’s overall product lifecycle by assisting in requirement analysis, workflow optimization, feature coordination, product monitoring, and operational support. The department helps bridge the gap between business requirements and technical implementation, ensuring that products are developed, tested, deployed, and maintained effectively. By supporting operational efficiency and product scalability, the team plays an important role in delivering reliable digital banking and credit solutions to clients and customers.

\section{Technical Infrastructure}

The OMS (Operations Management System) team at Hyperface utilizes a modern software development and cloud-based infrastructure environment to support fintech operations and product management activities. The department primarily works with web-based dashboards, API management platforms, cloud-hosted services, issue tracking systems, and collaboration tools for monitoring workflows and coordinating product development tasks. Development and testing activities are carried out using modern software technologies, version control systems, and continuous integration and deployment (CI/CD) pipelines.

The team also uses internal operational management platforms, analytics dashboards, database systems, and secure cloud infrastructure for transaction monitoring and workflow management. Collaboration and documentation tools are used extensively for requirement analysis, sprint planning, feature tracking, and communication between engineering, operations, and business teams. The infrastructure is designed to support secure, scalable, and high-availability fintech applications and digital banking services.

\section{Workflow and Methodologies}

The OMS team at Hyperface follows an Agile-based workflow for managing product development and operational activities. Work is generally organized into short development sprints where product requirements, feature updates, bug fixes, and operational improvements are planned and executed collaboratively. Daily stand-up meetings, sprint planning sessions, and review discussions are conducted to track progress and coordinate between product, engineering, and operations teams.

The department follows a simplified Software Development Life Cycle (SDLC) consisting of requirement gathering, feature planning, development, testing, deployment, and monitoring stages. Tasks and issues are managed using internal tracking tools and collaboration platforms. Since Hyperface operates as a fast-paced startup, the workflow is flexible and adaptive, allowing teams to quickly respond to changing business requirements, client feedback, and operational needs while maintaining product quality and system reliability.

%%%%%%%%%%%%%%%%%%%%%%%%%%%%%%%%%%%%%%%%%%%%%%%%%%%%%%%%%%%%%%%%%%%%%%%%%%%%%%%%%%%%%%%%%%%%%%%%%%%%
% NOTE: This section is meant to teach how to use "Acronyms" and "Glossaries" in the report. 
% 		So Remove or comment the following section and it's content entirely.
%%%%%%%%%%%%%%%%%%%%%%%%%%%%%%%%%%%%%%%%%%%%%%%%%%%%%%%%%%%%%%%%%%%%%%%%%%%%%%%%%%%%%%%%%%%%%%%%%%%%
